\section{Fractional Gaussian persistence}
\label{pr:section}


This section proves the probability estimates used in the divisor-box
arguments. All conditional probabilities below integrate the marks and
retain the gap or Gaussian environment. In particular, a fractional power
is taken \emph{after} any permitted mark window has been summed. The
conditioning and projection arguments preserve this convention.

Write $\mathbb P$ and $\mathbb E$ for probability and expectation. Comparison
constants are uniform on the compact parameter sets explicitly specified
below. They may depend on fixed curve amplitudes, endpoint boxes, density
margins, and regularity tolerances. None depends on the length of the
process, its allowed reserve, or a count in the stated central range.

\subsection{The Gaussian exponent and effective finite-time bounds}

Let $X,Y$ be independent stationary standard Ornstein--Uhlenbeck processes
with covariance $\mathbb E X_sX_t=e^{-|s-t|/2}$. For $0<\lambda,\rho<1$ put
\begin{align}
 A_t&=\{\sqrt\rho X_s+\sqrt{1-\rho}Y_s\ge0:0\le s\le t\},\notag\\
 q_t(X)&=\mathbb P_Y(A_t\mid X),\qquad
 F_{\lambda,\rho}(t)=\mathbb E_X q_t(X)^\lambda. \label{pr:defF}
\end{align}
We suppress the parameters when they are fixed. Finite-grid assertions
are passed to continuous paths by decreasing events on nested dense grids
and bounded convergence.

\begin{proposition}[Existence and effective bounds]\label{pr:existence}
The limit
\[
 \chi(\lambda,\rho)=\lim_{t\to\infty}-t^{-1}\log F_{\lambda,\rho}(t)
\]
exists. If $b(t)=-\log F(t)$ and
$p_g=(1+e^{-g/2})/(1-e^{-g/2})$, then, for $t,g>0$,
\begin{equation}
 \frac{b(t)}{p_g(t+g)}\le\chi(\lambda,\rho)
             \le\frac{b(t)}t,
 \qquad \frac\lambda2\le\chi(\lambda,\rho)\le\frac12.
 \label{pr:effective}
\end{equation}
Thus $g=2\log t$, $t>1$, gives an effective interval of width
$O(\log t/t)$. Finite Gaussian integrals give arbitrarily accurate
certified values of the endpoints of this interval.
\end{proposition}

\begin{proof}
The private survival indicators on two adjacent time intervals are
increasing functions of a Gaussian vector with nonnegative covariances.
Gaussian association \cite{Pitt}, in the finite form proved in
Lemma~\ref{pr:association}, first in $Y$ and then in $X$, gives
\begin{equation}
 F(s+t)\ge F(s)F(t).                                      \label{pr:super}
\end{equation}
The function $b$ is nonnegative, locally bounded and subadditive.
Writing $T=nt+r$, $0\le r<t$, proves that its limiting slope exists
and equals $\inf_{t>0}b(t)/t$.

For completeness, the Gaussian H\"older inequality used here is the
following finite-dimensional result: if a centered Gaussian block vector
has covariance bounded above by $\operatorname{diag}(p_i\Sigma_{ii})$,
then $\mathbb E\prod_i f_i(Z_i)\le\prod_i
(\mathbb E f_i(Z_i)^{p_i})^{1/p_i}$ for nonnegative functions.
This is \cite[Theorem~1(i)]{CDP}.
Consider $N$ intervals of length $t$, separated by gaps $g$. The OU Markov
property and conditional expectation at their endpoints imply that the
covariance of linear functionals on blocks $i,j$ is bounded in absolute
value by $r^{|i-j|}$ times their standard deviations, where $r=e^{-g/2}$.
Consequently, by $2ab\le a^2+b^2$,
\[
 \operatorname{Var}\Big(\sum_i L_i\Big)
 \le \left(1+2\sum_{k\ge1}r^k\right)\sum_i\operatorname{Var}(L_i)
 =p_g\sum_i\operatorname{Var}(L_i).
\]
This is covariance domination, uniformly in the grids and $N$.
Drop survival in the intervening gaps. Private Gaussian H\"older gives
$q_{\rm joint}\le\prod_i q_i^{1/p_g}$. Applying the same inequality to
the environmental expectation gives
\[
 F(Nt+(N-1)g)\le
 \mathbb E\prod_iq_i^{\lambda/p_g}\le F(t)^{N/p_g}.
\]
Only one factor $p_g$ remains. Taking $N\to\infty$ proves the first
inequality of \eqref{pr:effective}.

The composite process in \eqref{pr:defF} is itself a standard OU process.
Brownian time change and reflection give its ordinary survival probability
\begin{equation}
 a(t)=\pi^{-1}\arcsin(e^{-t/2}).                            \label{pr:ordinaryOU}
\end{equation}
The inequalities $q^\lambda\ge q$ and Jensen give
$a(t)\le F(t)\le a(t)^\lambda$, proving the range of $\chi$ and also
$b(t)\le t/2+\log\pi$. This last bound makes the width assertion explicit.

Here are the finite-grid error bounds, to specify the claimed effective
meaning. At mesh size $\Delta=t/n$, let $G_h$ denote the fractional grid
survival moment with common threshold $h$. Conditional on two composite
endpoints $x,y\ge h$, the intervening crossing probability is
$\exp[-xy/\sinh(\Delta/2)]$. Therefore
\begin{align}
 F(t)&\le G_0,\notag\\
 F(t)&\ge G_h-
       \{n\exp[-h^2/\sinh(\Delta/2)]\}^{\lambda}.            \label{pr:griderror}
\end{align}
Indeed the conditional error probability is first raised to $\lambda$;
Jensen bounds its expectation by its annealed probability to that power.
Translation of the composite Gaussian grid gives
\begin{equation}
 |G_h-G_0|\le
 \left\{\frac{|h|}{\sqrt{2\pi}}
             \sqrt{1+n\tanh(\Delta/4)}\right\}^{\lambda}.
 \label{pr:gridthreshold}
\end{equation}
The squared translation norm is
$h^2\boldsymbol1^T\Sigma^{-1}\boldsymbol1
=h^2[1+n\tanh(\Delta/4)]$, where
$\Sigma_{ij}=e^{-|i-j|\Delta/2}$. Equal-covariance Gaussian total
variation is at most that norm divided by $\sqrt{2\pi}$.
Taking, for example, $h=\Delta^{1/4}$ makes both errors vanish.
On a bounded cube the Gaussian density is uniformly continuous and the
orthant integrand is bounded and monotone; displaced-threshold upper and
lower cube sums, followed by Gaussian tail bounds, give certified finite
quadratures. This completes the effective construction without any
assumption about a divisor-volume asymptotic.
\end{proof}

\subsection{Tilted environments and the private endpoint filter}

Write $\nu_{t,\lambda}$ for the environmental probability measure with
density $q_t^\lambda/F(t)$ relative to the OU law. We require endpoint
estimates under this measure, including the private endpoint after it
has been conditioned to survive against the actual environment.

\begin{lemma}[Order of the tilted laws]\label{pr:order}
For $0\le\lambda\le1$,
\begin{equation}
 \nu_{t,0}\preceq_{\rm st}\nu_{t,\lambda}
                         \preceq_{\rm st}\nu_{t,1}.          \label{pr:storder}
\end{equation}
For each $a<1/2$,
\begin{equation}
 \sup_{t\ge0,\,0\le u\le t}
      \mathbb E_{\nu_{t,\lambda}}e^{aX_u^2}<\infty.          \label{pr:envmoment}
\end{equation}
A nonnegative decreasing environmental path functional has tilted
expectation at most its ordinary expectation.
\end{lemma}

\begin{proof}
On a finite grid the OU precision matrix has nonpositive off-diagonal
entries, so its density is multivariate totally positive of order two
(MTP$_2$). Set $Z=-Y$ and $a_\rho=\sqrt{\rho/(1-\rho)}$.
The set $\{z_i\le a_\rho x_i\ \forall i\}$ is closed under coordinatewise
minimum and maximum. Its indicator is MTP$_2$. Multiplication by the
MTP$_2$ Gaussian density and marginalization over $Z$ show that $q(x)$
is MTP$_2$. One can insert an independent product Gaussian density in
$x$ to ensure integrability; its factors cancel in the defining inequality.
Thus $\varphi_\Sigma(x)q(x)^\lambda$ is MTP$_2$ and is associated.
The closure and association properties used here are given in
\cite[Proposition~3.4 and Section~3]{Fallat}.
Association under this tilted measure, applied to an increasing function
and $q^{1-\lambda}$, proves the right comparison in
\eqref{pr:storder}. Ordinary Gaussian association proves the left one.
The decreasing-functional assertion is the same left comparison.

Under $\nu_{t,1}$ rotate to
$Z=\sqrt\rho X+\sqrt{1-\rho}Y$ and an independent standard OU process
$V$. Only $Z$ is conditioned to be positive. Its density at time $u$ is
\[
 \frac{\varphi(z)\mathbb P_z(\tau_0>u)
                    \mathbb P_z(\tau_0>t-u)}{a(t)}
                 \boldsymbol1_{\{z>0\}},
 \qquad
 \mathbb P_z(\tau_0>v)=2\Phi\left(\frac z{\sqrt{e^v-1}}\right)-1.
\]
For each side of length at least one the latter probability is at most
$Cz e^{-v/2}$; a shorter side contributes at most one. Separating the
cases with zero, one or two long sides and using \eqref{pr:ordinaryOU}
bounds the density by $C(1+z)^2\varphi(z)$, uniformly in $u,t$.
The rotation gives uniform Gaussian exponential moments for $X_u$ under
$\nu_{t,1}$. Apply \eqref{pr:storder} to the increasing function
$e^{a(x_+)^2}$ and the decreasing function $e^{a(x_-)^2}$ separately.
This proves \eqref{pr:envmoment}; dense grids give the path assertions.
\end{proof}

\begin{lemma}[Uniform private filter estimate]\label{pr:filter}
On compact subsets of $0<\lambda,\rho<1$ there exist $a>0$ and $C<\infty$
such that
\begin{equation}
 \mathbb E_X\left[e^{aX_t^2}
  \left(\mathbb E_Y[\boldsymbol1_{A_t}e^{aY_t^2}\mid X]\right)^\lambda
             \right]\le C F(t).                            \label{pr:mixed}
\end{equation}
The same assertion holds with simultaneous Gaussian weights at both
initial and terminal endpoints, after decreasing $a$.
\end{lemma}

\begin{proof}
Fix the environmental path and let the private wall be
$Y_s\ge-a_\rho X_s$. The filtered endpoint density has the form
$\mu_v(y)\propto\varphi(y)g_v(y)$, where $g_v$ is increasing and
log-concave. Indeed the conditional OU bridge mean is increasing in its
endpoint, and its constraints are convex; marginal log-concavity is
Pr\'ekopa's marginalization theorem
\cite[Theorem~6, p.~342]{Prekopa1973}, first on a finite grid and then by
a limit.
Thus $\mu_v$ is $1$-strongly log-concave and its mode $m_v$ is nonnegative.
For any such density with mode $m$,
\begin{equation}
 |Z-m|\preceq_{\rm st}|N(0,1)|,
 \qquad |\mathbb E Z-m|\le\sqrt{2/\pi}.                     \label{pr:mode-tail}
\end{equation}
To verify this directly, both functions
$\mu(m+r)e^{r^2/2}$ and $\mu(m-r)e^{r^2/2}$ decrease for $r\ge0$.
Their sum therefore gives likelihood-ratio domination of the radial
half-normal density after normalization.
The density is also dominated in monotone likelihood-ratio order by
$m+|N|$: on $y\ge m$ its ratio to $e^{-(y-m)^2/2}$ decreases, while
on $y<m$ the comparison density is zero. OU transition kernels are TP$_2$;
composition with killing indicators and integration preserve this property.
The likelihood-ratio convention allowing zero densities agrees with
\cite[Theorem~1 and Remark~16]{DuembgenMoesching2023}; kernel composition
here follows from the MTP$_2$ marginalization just used. Consequently
this order propagates through a killed OU segment.
Raising a wall also increases its normalized endpoint law, by association
under the original truncated MTP$_2$ path measure.

Consider a segment of length $v\in[1,2]$, and put $r=e^{-v/2}$,
$B=\max(0,\sup_{\rm segment}(-a_\rho X))$. Raise the incoming law to
$m+|N|$ and the wall to $B$. If $rm\ge B+2$, an OU-noise tube of radius
one ensures survival with a fixed positive probability. The unconditional
endpoint above $rm$ has a uniform Gaussian tail, so its conditional mean
is at most $rm+C$. If $rm<B+2$, restrict the initial value to
$[e(B+3),e(B+3)+1]$ and use the same noise tube. The restriction has
probability at least $c e^{-C(B+1)^2}$. Dividing the unconditional
Gaussian endpoint tail by this probability and integrating yields a
conditional mean at most $rm+C(B+1)$. Equation \eqref{pr:mode-tail}
therefore gives
\begin{equation}
 m_{\rm new}\le e^{-v/2}m+C(B+1).                            \label{pr:mode-recursion}
\end{equation}
An initial segment of length at most one is handled only once: raise its
wall to its supremum and retain a fixed initial interval above
$e(B+3)$. The same tube and tail argument bounds its mode by $C(B+1)$.
There is no iteration with a contraction tending to one.

Partition backwards from the terminal time into intervals $I_j$ of length
at most two. Iterating \eqref{pr:mode-recursion}, including the initial
truncation, gives
\begin{equation}
 m_t(X)\le C\left(1+\sum_{j\ge0}e^{-j/2}
                         \sup_{s\in I_j}(-a_\rho X_s)_+\right)=:\mathcal M_t(X).
 \label{pr:mode-majorant}
\end{equation}
OU suprema on bounded intervals have uniform Gaussian tails. Weighted
Cauchy--Schwarz and Jensen imply
$\sup_t\mathbb E e^{c\mathcal M_t^2}<\infty$ for some $c>0$; independence
of the intervals is not needed. The majorant is decreasing in $X$, so
Lemma~\ref{pr:order} transfers this bound to $\nu_{t,\lambda}$.
Finally, for $a$ small, \eqref{pr:mode-tail} gives
\[
 \mathbb E[e^{aY_t^2}\mid A_t,X]
 \le (1-4a)^{-1/2}e^{2am_t(X)^2}.
\]
Multiply by $q_t^\lambda$, use \eqref{pr:envmoment} and the preceding
mode bound, and apply Cauchy--Schwarz under $\nu_{t,\lambda}$.
This proves \eqref{pr:mixed}. Reversibility treats the other endpoint;
Cauchy--Schwarz first inside the private law and then in the environmental
law treats both endpoints simultaneously. The private law was never
replaced by a fresh Gaussian after survival was imposed.
\end{proof}

\begin{theorem}[Constant-factor Gaussian survival]\label{pr:gaussian}
Uniformly on compact subsets of $0<\lambda,\rho<1$,
\begin{equation}
 F(s)F(t)\le F(s+t)\le C F(s)F(t),\qquad
 C^{-1}e^{-\chi t}\le F(t)\le e^{-\chi t}.                  \label{pr:quasimult}
\end{equation}
The second comparison is valid for $t\ge1$. There exist compatible,
sufficiently wide fixed positive endpoint boxes for both Gaussian
channels that can be imposed at a constant cost. These are the boxes
constructed below, or fixed boxes containing them; this assertion does
not prescribe an arbitrary narrow endpoint box.
\end{theorem}

\begin{proof}
The Mehler density relative to two independent standard Gaussians is
\[
 k_r(x,y)=\frac1{\sqrt{1-r^2}}
 \exp\left(\frac{2rxy-r^2(x^2+y^2)}{2(1-r^2)}\right)
 \le\frac1{\sqrt{1-r^2}}
       \exp\left(\frac{r(x^2+y^2)}{2(1+r)}\right).
\]
Choose a fixed gap $g$ so large that $r=e^{-g/2}$ makes the exponent
smaller than the coefficient allowed in Lemma~\ref{pr:filter}.
Drop survival across the gap. Apply this density bound inside the private
conditional integral and then to the environmental transition. The two
reference blocks are independent; \eqref{pr:mixed} at their adjacent
endpoints gives $F(s+g+t)\le C F(s)F(t)$. On the other hand,
\eqref{pr:super} gives $F(s+g+t)\ge F(s+t)F(g)$. This proves the first
comparison. Iteration at $nt$ and Proposition~\ref{pr:existence} yield
the second one.

By \eqref{pr:mixed}, excluding any of the four Gaussian endpoints from
$[-R,R]$ costs at most $Ce^{-cR^2}F(t)$ in fractional mass. Choose fixed
$R$ so that their union costs at most $F(t)/2$. Add to each channel
$A(e^{-s/2}+e^{-(t-s)/2})$, with any fixed $A\ge2R$. This preserves
survival and moves all retained endpoints into the fixed positive box
$[A-R,A(1+e^{-1/2})+R]$, uniformly for $t\ge1$.
The Cameron--Martin translation formula \cite{CameronMartin1944}, in
the form including an independent Gaussian initial value given in
\cite[Theorem~2.1]{BarbourRossZheng2024}, applies after the OU-to-Brownian
time change. The density of this shift depends only on those endpoints and is bounded above and below
on the box, uniformly for $t\ge1$. Change variables first in the private
integral and then in the environmental integral. The positive-box event
therefore retains a fixed fraction of $F(t)$. A required fixed positive
lower endpoint level can be accommodated by choosing $A$ first and then
retaining this sufficiently wide box (or a fixed containing box).
No restriction to an arbitrary prescribed narrower box is inferred.
\end{proof}

\subsection{Parameter stability, integrable walls, and Brownian reserves}

\begin{proposition}[Local Lipschitz continuity]\label{pr:lipschitz}
On each compact subset of $(0,1)^2$,
\[
 |\chi(\lambda,\rho)-\chi(\lambda',\rho')|
             \le C(|\lambda-\lambda'|+|\rho-\rho'|).
\]
\end{proposition}
\begin{proof}
Whiten a finite Gaussian grid. Its survival indicator is a fixed function
$f$, and $q_\rho=T_s f$, where $T_s$ is the Gaussian noise semigroup and
$e^{-s}=\sqrt\rho$. In any unit direction Gaussian integration by parts
and the one-dimensional rearrangement inequality give
\[
 |\nabla q_\rho|\le\sqrt{\frac\rho{1-\rho}} I(q_\rho),
 \qquad I(v)=\varphi(\Phi^{-1}v).
\]
For clarity, the rearrangement step fixes the value $v=\mathbb E f$ and
maximizes $|\mathbb E[Z_1 f(Z)]|$ over $0\le f\le1$; the maximizer is a
half-space in the first coordinate, with value $I(v)$. Gaussian tail
bounds give $I(v)^2\le Cv^2(1-\log v)$, including both endpoints.
Differentiating under the Gaussian integral and integrating by parts,
\[
 \partial_s\log\mathbb E q_s^\lambda
 =\lambda(1-\lambda)
   \frac{\mathbb E q_s^{\lambda-2}|\nabla q_s|^2}
        {\mathbb E q_s^\lambda}.
\]
Convexity of $K(\lambda)=\log\mathbb E q^\lambda$, with $K(0)=0$, gives
$\mathbb E_{q^\lambda}[-\log q]\le-K(\lambda)/\lambda$.
Since $\mathbb E q^\lambda\ge a(t)\ge\pi^{-1}e^{-t/2}$, both the
$\lambda$ derivative and, on a compact $\rho$ interval, the $\rho$
derivative are bounded in absolute value by $C(t+1)$, independently of
the grid dimension. Smooth approximation justifies the derivatives when
$f$ is an indicator. Integrating these bounds, passing to dense grids,
dividing by $t$, and taking $t\to\infty$ proves the assertion.
\end{proof}

For a piecewise smooth shift $h$ on $[0,t]$, define
\begin{align*}
 \|h\|_{\mathcal H}^2&=h(0)^2+\int_0^t(h'+h/2)^2\,ds,\\
 J_t(h)&=|h(0)/2-h'(0)|+|h'(t)+h(t)/2|
                         +\int_0^t|h/4-h''|\,ds.
\end{align*}
Derivative jumps are included in the last total variation.

\begin{lemma}[Bounded-action walls]\label{pr:wall}
Replacing the zero composite OU wall by a deterministic wall $b$ changes
$F(t)$ by a factor bounded above and below whenever
$J_t(b/\sqrt\rho)+\|b/\sqrt\rho\|_{\mathcal H}^2$ is bounded.
An endpoint-restricted version holds with compatible sufficiently wide
boxes chosen after the wall bound, or with the endpoint boxes supplied
by Theorem~\ref{pr:gaussian} translated together with the path.
It is not a statement for arbitrary prescribed unchanged endpoint boxes.
\end{lemma}
\begin{proof}
The Cameron--Martin linear functional is
\[
 L_h(X)=(h(0)/2-h'(0))X_0+(h'(t)+h(t)/2)X_t
                   +\int_0^t(h/4-h'')X_s\,ds.
\]
Environmental translation alone gives exactly
$F_b(t)/F(t)=\mathbb E_{\nu_{t,\lambda}}
\exp[-L_h(X)-\|h\|_{\mathcal H}^2/2]$, with $h=b/\sqrt\rho$ and the
consistent translation sign. The absolute variation of the coefficients
of $L_h$ is bounded. Jensen with its normalized absolute coefficient
measure and \eqref{pr:envmoment} bounds both exponential signs uniformly.
Jensen also supplies a positive lower bound, or equivalently use the
bound for $\mathbb E_{\nu}|L_h|$. Mixed endpoint moments, followed by
Cauchy--Schwarz and truncation, give the same conclusion for compatible
sufficiently wide boxes chosen after the fixed wall bound. Equivalently,
translate the endpoint
boxes in the same environmental change of variables. A fixed feasible
interior margin and a finite initial or terminal tube permit the
sufficiently wide compatible central boxes used below; their constants
are chosen after the wall and drift bounds.
This argument does not change the power $\lambda$.
\end{proof}

\begin{corollary}[Brownian one-end estimate]\label{pr:brownian}
Let independent Brownian channels have environmental variance $\sigma_E^2$
and private variance $\sigma_C^2$, bounded above and below by positive
constants. Put $\rho=\sigma_E^2/(\sigma_E^2+\sigma_C^2)$. From reserve
$h\ge1$, over duration $n$, with $|\delta|\sqrt n\le D$, the conditional
private survival probability $Q$ satisfies
\begin{equation}
 \mathbb E Q^\lambda\asymp
 B_n(h):=\min\{1,((1+h)/\sqrt n)^{2\chi(\lambda,\rho)}\}.
 \label{pr:Bn}
\end{equation}
The upper bound allows a fixed soft wall $-h-gs^\gamma$, $0<\gamma<1/2$.
For a sufficiently large fixed minimum reserve, the lower bound permits
a strong wall $-h+gs^\gamma$, positive central terminal risk, and
compatible sufficiently wide central terminal boxes of the separate
channels. Choose these boxes after the compact-law, wall, and drift
bounds; their locations and widths are then fixed uniformly in $n,h$.
\end{corollary}
\begin{proof}
For $h^2<n$, scale from time $h^2$ and use Lamperti time $u=\log(s/h^2)$.
Its length is $\log(n/h^2)$. Reserve, power wall and bounded scaled drift
become bounded multiples of
\[
 e^{-u/2},\qquad e^{-(1/2-\gamma)u},\qquad
 e^{-(\log(n/h^2)-u)/2}.
\]
Their $J$ and Hilbert norms are bounded. Drop the initial time interval
for the upper bound and apply Lemma~\ref{pr:wall} and
Theorem~\ref{pr:gaussian}. For the lower bound prepare, on a fixed multiple
of $h^2$, both channels in positive central boxes using a fixed Brownian
tube. Brownian scaling makes its probability and its endpoint transition
density bounds uniform. Thereafter use the endpoint version of those
lemmas. The preparation has fixed cost before the fractional power as
well as after it. A large fixed minimum reserve makes the normalized
strong-wall amplitude as small as required. If $h$ is comparable to or
larger than $\sqrt n$, a fixed-scale positive tube and monotonicity give
a constant lower bound; the upper bound is one. Reflection of both
channels proves the same assertions for the opposite risk sign.
\end{proof}

\subsection{Independent exponential gaps: relative coupling}

Let $C_i$ be iid marks in a fixed finite alphabet, with probabilities in
a compact subset of its open simplex. Their weights have two values,
\[
 w(C_i)=l-bB_i\in[w_*,l],\qquad
 p=\mathbb P(B_i=1),\qquad \omega=l-bp,
\]
where all parameters range over compact sets, $b,w_*>0$, and $p$ stays
away from zero and one. Independent gaps $E_i$ have exponential mean
$\tau$ in a positive compact interval. Write
\begin{equation}
 S_j=\sum_{i\le j}(w(C_i)-E_i),\quad
 \delta=\omega-\tau,\quad
 \rho_\tau=\frac{\tau^2}{\tau^2+b^2p(1-p)},\quad
 F_n(h)=\mathbb E\mathbb P_C\{\min_{j\le n}S_j\ge-h\mid E\}^{\lambda}.
 \label{pr:walk}
\end{equation}
The following arguments also apply to $-S$, with all monotonicities
reversed together. Assume initially $|\delta|\sqrt n\le D$.

\begin{lemma}[Crude comparisons before the exponent is known]
\label{pr:crude}
Uniformly for $h\ge1$, $u\ge0$,
\begin{equation}
 F_n(h)\ge c_D n^{-1/2},\qquad
 F_n(h+u)\le C_D(1+u)F_n(h),\qquad
 F_{2n}(h)\ge c_D F_n(h).                                 \label{pr:crudeineq}
\end{equation}
The drift bound for a doubled interval is enlarged by a fixed factor.
The assertions remain valid for a fixed soft sub-square-root power wall,
with a fixed polynomial in $1+u$ allowed in the middle inequality.
\end{lemma}
\begin{proof}
We detail the elementary annealed preparation used here. For the centered
walk stopped on leaving $[-h_0,B]$, upper overshoots are bounded and lower
overshoots have exponential tails; for the reversed walk these roles
are interchanged. Optional stopping of the walk gives an upper-exit
probability comparable to $1/B$, from a fixed reserve $h_0>0$.
Stopping its square minus the variance times time bounds the mean exit
time by $C B$: on an upper exit its squared displacement is $O(B^2)$
and that exit has probability $O(1/B)$; on a lower exit the second
moment is bounded. Truncation first justifies both stopping identities.
By Markov's inequality, a sufficiently large fixed multiple of $B^2$
retains a fixed fraction of the upper exits. From height $B$, a positive
scaled tube over a further comparable time has probability bounded below
by the functional central limit theorem, uniformly on our compact laws.
It can end above any prescribed fixed multiple of $B$.
The lower-overshoot memoryless law and bounded upper overshoot provide
uniformity at entrance. The stopped centered gap and mark martingales
have second moments at most $C\mathbb E\tau_{\rm exit}\le CB$.
Their probability of exceeding $RB$ before the retained entrance is
therefore at most $C/(R^2B)$. Choosing fixed $R$ large retains a fixed
fraction of the $B^{-1}$ entrance probability with both channel sums
central; the subsequent tube preserves centrality at the final time.
Small drift is handled by changing the exponential
mean on the retained central gap-sum tube: if $m=O(B^2)$ and
$|\delta|\sqrt m=O_D(1)$, its log density ratio is $O_D(1)$ there.
The same proof, or its sign reversal, therefore gives a preparation of
length $O((1+u)^2)$ from fixed reserve to height at least $u$, of
annealed probability at least $c_D/(1+u)$.

For $u\le\sqrt n$, let $p_u$ be its conditional private probability.
Concatenating with an independent $n$-step survival from $h+u$, and
using $\mathbb E p_u^\lambda\ge\mathbb E p_u$, gives
$F_{n+m}(h)\ge c_D(1+u)^{-1}F_n(h+u)$. Time monotonicity gives the
middle inequality. The same preparation and tube give ordinary survival
at least $c_D/\sqrt n$, hence the first inequality. For $u>\sqrt n$
these two inequalities and $F_n(h+u)\le1$ suffice.

For doubling, remove environments whose centered prefix gap sum is
unfavorable by more than $R\sqrt n$. Survival is monotone in the opposite
direction to this unfavorable event, so product association bounds its
$F_n$-weighted mass by its unweighted probability times $F_n$.
Choose $R$ large. On the retained environments an appropriate favorable
private centered endpoint of size $(R+D+c)\sqrt n$ has a fixed positive
probability. Its mark event is monotone in the same direction as survival,
so their intersection retains that fraction of the conditional survival
probability. Its ending reserve is at least $c\sqrt n$.
An independent further $n$ steps survive from this height with annealed
probability at least a constant. If $f$ is the prefix mark probability
and $Z$ its completed probability as a function of the future environment,
then $0\le Z\le f$ and $\mathbb E Z\ge cf$, whence
$\mathbb E Z^\lambda\ge cf^\lambda$. Integration proves doubling.
All finite small lengths are absorbed using compactness and a fixed
positive gap/mark event. For a soft wall use
$(s+t)^\gamma\le s^\gamma+t^\gamma$ after a prefix; its additional reserve
is polynomial in the prefix length. The preparation and favorable
continuation remain sufficient. No value of $\chi$ is used in this proof.
\end{proof}

\begin{lemma}[Relative cost of a bad block]\label{pr:relativebad}
Let $\mathcal D$ depend on the first $b\le n/2$ slots, including any
coupling auxiliaries. Then
\begin{equation}
 \mathbb E\mathbb P_C(\text{survival},\mathcal D\mid
                  \text{extended environment})^\lambda
 \le C_D(1+b)^2\mathbb P(\mathcal D)^{\lambda/4}F_n(h).
 \label{pr:earlybad}
\end{equation}
If instead $\mathcal D$ begins after $n/4$, its relative cost is at most
$C_D\mathbb P(\mathcal D)^\lambda F_n(h)$. Fixed fractions in these
statements may be changed, with corresponding constants. The same
bounds hold for the Gaussian comparison walk and for both risk signs.
\end{lemma}
\begin{proof}
If positive jumps are bounded, discard prefix survival and use the
endpoint bound $S_b\le lb$. Independence of the remainder, Jensen over
the prefix, and Lemma~\ref{pr:crude} give the stronger upper bound
$C(1+b)\mathbb P(\mathcal D)^\lambda F_n(h)$.
For unbounded positive jumps divide positive prefix endpoints into unit
bins $[k,k+1)$. Their probabilities are bounded both by
$\mathbb P(\mathcal D)$ and by
$C\exp[-c\min(k^2/b,k)]$, with the fixed $O_D(\sqrt b)$ drift displacement
absorbed in the constants and a first group of bins. Apply the reserve
comparison bin by bin. Using the geometric mean of the two probability
bounds and summing the Gaussian/exponential tail proves
\eqref{pr:earlybad}; the displayed exponent $\lambda/4$ and quadratic
polynomial are deliberately weaker than necessary. Negative endpoints
can only decrease the available reserve. The Gaussian case has the same
tail estimate. For a late bad block retain necessary survival on the
first $\lfloor n/8\rfloor$ slots, independent of that block. Jensen and
finitely many doublings prove the stated bound. These are relative
bounds, not an assertion that the bad event is independent of survival.
\end{proof}

\begin{lemma}[Relative control of the initial coupling block]
\label{pr:initial-coupling-block}
Fix the compact laws and the bounded scaled-drift range in
Lemma~\ref{pr:relativebad}, and let the fractional exponent satisfy
$\lambda\in[\lambda_*,1]$, where $\lambda_*>0$ is fixed. Fix an integer
$b\ge1$. Couple the centered
exponential and binary channels separately to their scalar Gaussian
channels on the first $b$ slots, independently of all later blocks.
The environmental coupling is outside the private conditional
probability, and the private coupling is inside it.
For every prescribed $\varepsilon>0$ there is a finite $H_b$, depending
only on these fixed data and $\varepsilon$, such that, whenever
$n\ge2b$ and $h\ge H_b$, the following assertion holds for each
marginal survival functional covered by Lemma~\ref{pr:relativebad}.
There is an initial good event $\mathcal G_b$ on which the discrepancy
of the two initial risk paths is at most $H_b/8$, and
\[
 \mathbb E\bigl[
   \mathbb P_{\mathrm{private}}
     (\mathrm{survival},\mathcal G_b^c
           \mid\mathrm{extended\ environment})^{\lambda}\bigr]
       \le \varepsilon F_n(h).
\]
The same $H_b$ works for both risk signs. A prescribed fixed strong or
soft curve on this initial finite interval may also be absorbed in
$H_b/8$ by increasing $H_b$ before lengths vary.
\end{lemma}
\begin{proof}
Let $M_b$ be the maximum of the absolute centered partial sums of
both original channels on the first $b$ slots and of the absolute
Gaussian suprema on the corresponding time interval. Enlarge its
definition by the fixed compact coefficients needed to reconstruct
the risk. No independence between the two members of a channel
coupling is required for the following union bound. Exponential
moments of the gaps, boundedness of the binary increments, and the
Gaussian maximal tail give constants $C_b,c_b>0$ such that
\[
 \mathbb P\{M_b>v\}\le C_b e^{-c_b v}
 \qquad(v\ge C_b).
\]
The constants are independent of $n$ and $h$. The deterministic drift
on the first $b$ slots is bounded in terms of the fixed data, because
$n\ge2b$ and the full-horizon scaled drift is bounded.
Choose a fixed compact-data constant $C_0$ so that, on
$\mathcal G_b=\{M_b\le H_b/C_0\}$, the two initial risk paths differ
by at most $H_b/8$; make the same choice large enough to pay the
initial drift. A fixed curve on this finite interval has a finite
maximum, which is paid by the final increase of $H_b$.

The bad event reads only the first $b$ slots and their coupling
auxiliaries. Lemma~\ref{pr:relativebad} therefore gives
\[
 \mathbb E\bigl[
   \mathbb P_{\mathrm{private}}
      (\mathrm{survival},\mathcal G_b^c
           \mid\mathrm{extended\ environment})^{\lambda}\bigr]
 \le C(1+b)^2
       \bigl(C_b e^{-c_bH_b/C_0}\bigr)^{\lambda/4}F_n(h).
\]
After $b$ has been fixed, choose $H_b$ so that its coefficient is at
most $\varepsilon$, uniformly for $\lambda\ge\lambda_*$. This choice
does not depend on a persistence exponent or on the eventual reserve
$h$. Since $h\ge H_b$, a change by $H_b/8$ lies between the reserves
$7h/8$ and $9h/8$. The argument bounds the initial error only on
$\mathcal G_b$; it has assigned no deterministic bound to any
exponential gap. The finitely many experiments $n<2b$ are outside this
lemma and are handled by their explicitly supported finite cases.
\end{proof}

\begin{theorem}[Reference one-end theorem]\label{pr:one}
For $|\delta|\sqrt n\le D$, $h\ge1$, and either sign of $S$,
\begin{equation}
 F_n(h)\asymp_D
 \min\{1,((1+h)/\sqrt n)^{2\chi(\lambda,\rho_\tau)}\}.
 \label{pr:ONE}
\end{equation}
The upper estimate allows a fixed soft power wall of exponent below
$1/2$. For $h\ge H_0$ the lower estimate may retain a fixed strong curve,
a fixed positive central unreserved endpoint, and central gap, binary,
and full-color endpoints. The central boxes are chosen compatible and
sufficiently wide after the compact laws, wall, $D$, and $D'$ are fixed,
before the coupling accuracy and minimum reserve are chosen.
For this endpoint version take $h\le D'\sqrt n$;
larger reserves are not needed for the endpoint assertion.
\end{theorem}
\begin{proof}
We use a finite strong approximation with explicit failure tails, not an
unconditional limiting invariance principle. The finite scalar coupling of
Sakhanenko \cite{Sakhanenko1984}, in the finite iid form of
Waudby-Smith, Larsson, and Ramdas \cite[Lemma~4.2, p.~12]{WLR}, followed
by Markov's inequality (as in their equation~(17)), gives
\[
 \mathbb P\!\left\{\max_{j\le L}
       \left|\sum_{i\le j}(Z_i-G_i)\right|\ge z\right\}
 \le (1+s_*\sqrt{Lv})e^{-cs_*z},
\]
where $v=\operatorname{Var}(Z)$, $c>0$ is universal, and $s_*$ is strictly
below the Sakhanenko parameter. We check this uniformly for the two
centered channels $Z=E-\tau$ and $Z=w(C)-\omega$ before conditioning.
Write $\tau\in[\tau_-,\tau_+]$, $b\in[b_-,b_+]$, and
$p\in[p_-,1-p_-]$ for their fixed compact bounds. Set
\[
 v_- =\min\{\tau_-^2,b_-^2p_-(1-p_-)\}>0,\quad
 v_+ =\max\{\tau_+^2,b_+^2/4\},\quad t_0=(2\tau_+)^{-1}.
\]
For $E=\tau U$ with $U$ exponential of mean one,
\[
 \mathbb E|E-\tau|^3e^{t_0|E-\tau|}
 \le e^{1/2}\tau_+^3\int_0^\infty(u+1)^3e^{-u/2}\,du
 =158e^{1/2}\tau_+^3.
\]
For the binary channel $|Z|\le b_+$, so the corresponding moment is at
most $b_+^3e^{t_0b_+}$. Let $M$ be the larger of these two bounds,
$t_1=\min\{t_0,v_-/(2M)\}$, and $s_*=t_1/2$. Then
$t_1\mathbb E|Z|^3e^{t_1|Z|}\le v_-/2\le v$; Definition~2.1 and
(6) of \cite{WLR} give a Sakhanenko parameter at least $t_1>s_*$.
These estimates also hold for $-Z$. For any fixed $A>0$ and $L\ge2$,
choose $C_A$ with $cs_*C_A>A+1/2$ and take $z=C_A\log L$. The displayed
tail is at most $(1+s_*\sqrt{v_+})L^{1/2-cs_*C_A}\le C L^{-A}$.
Increasing $C_A$ if needed gives error $C_A\log L$ with failure at most
$C_A L^{-A}$, uniformly in the compact laws. The bounded initial block
is treated separately below; no conditioned coupling is asserted.

Couple gaps and private binary marks separately, using independent
couplings on successive dyadic blocks, starting at a fixed scale $L_0$.
Include Brownian interpolation oscillations in the bad events; their
Gaussian bounds have the same arbitrarily high polynomial decay after
increasing $C_A$. A block intersecting the horizon either ends by $n/2$
or begins after $n/4$, apart from a bounded initial case. Apply
Lemma~\ref{pr:relativebad} on both marginal processes. The total
fractional failure mass is at most
\begin{equation}
 C F_n(h)\sum_{L\ge L_0,\ L\ {\rm dyadic}}L^{2-A\lambda/4},
 \label{pr:KMTrelative}
\end{equation}
and the analogous bound holds with Brownian survival in place of $F_n$.
Choose $A$ first, and then $L_0$, so this is any prescribed small
fraction of the corresponding marginal survival moment.

First fix a bounded-action envelope and the relative reserve interval
$[7h/8,9h/8]$. Fix its Gaussian and compatible endpoint lower constants.
Choose the coupling tail exponent and initial dyadic scale so the later
relative error sum is below the required fraction of these constants.
Apply Lemma~\ref{pr:initial-coupling-block} at that fixed initial scale;
choose its minimum reserve last. On the resulting good event the
initial discrepancy is at most $h/8$ and the later discrepancy is at most
$C\log(1+s/L_0)+C\log^2(1+s/L_0)$. After Lamperti transformation the latter
is a bounded-action curve of Lemma~\ref{pr:wall}. Increasing the minimum
reserve reduces its normalized amplitude; it does not enlarge the
chosen relative reserve interval. The initial bad event is paid at the
same marginal survival scale, not by a deterministic bound on gaps.
Finally choose the bounded-horizon threshold. None of these choices
uses the discrete normalization being proved.

The coupling is used on an extended probability space: the coupled gap
and Brownian environment is outside the private conditional probability;
the independent mark/Brownian coupling is inside it. The two marginal
conditional probabilities consequently remain the genuine $q_D(E)$ and
$q_B(W)$. On good paths the deterministic inclusions give
\[
 q_D\le q_{B,\mathrm{soft}}+\Delta_D,
 \qquad q_{B,\mathrm{hard}}\le q_D+\Delta_B.
\]
Subadditivity and \eqref{pr:KMTrelative} absorb the first error against
the discrete moment and the second against the Brownian moment.
Corollary~\ref{pr:brownian} proves \eqref{pr:ONE}, without presupposing
its normalization on either side. Reserves below the fixed threshold
are handled by the preparation and reserve comparisons of
Lemma~\ref{pr:crude}. The same argument retains Brownian strong curves
and endpoint boxes; its errors are still bounded by the flat relative
majorant.

Only the binary mark was coupled. Conditional on the entire binary path,
the full colors in each category remain independent with their conditional
probabilities. If the binary endpoint is central, their means differ
from the unconditional color means by $O(\sqrt n)$. Enlarging a fixed
central color window once, the multinomial variance bound makes all
colors central with conditional probability at least $3/4$, uniformly
for every retained binary path. This is a constant loss before the
fractional power. It requires no unstated multivariate KMT theorem.
\end{proof}

\subsection{Exact Gamma conditioning and relative local regularity}

Let $T_n=\sum_{i=1}^nE_i$ and condition on $T_n=A>0$, where
$|A-n\tau|\le D\sqrt n$ and $n\ge1$. Write
$0<\tau_-\le\tau\le\tau_+<\infty$ for the fixed gap-mean range.
All such conditioning is by a regular conditional
density. Write $\mathbb E_A$ for the resulting environmental expectation.
The private marks remain independent of the gaps and of one another.

\begin{lemma}[Selected-gap density]\label{pr:gamma}
For a deterministic proper set $I$ of $m<n$ coordinates, the conditional density
relative to independent mean-$\tau$ exponentials is exactly
\begin{equation}
 \frac{f_{n-m}(A-\sum_{i\in I}E_i)}{f_n(A)},\qquad
 f_k(x)=\frac{x^{k-1}e^{-x/\tau}}{\tau^k\Gamma(k)}
                          \boldsymbol1_{\{x>0\}}.          \label{pr:RN}
\end{equation}
There is a fixed threshold
$N_\Gamma=N_\Gamma(D,\tau_-,\tau_+)\ge
\max\{1,(2D/\tau_-)^2\}$ such that the following two upper bounds
hold for $n\ge N_\Gamma$: if $m\le cn$, $c<1$, the ratio is bounded
above uniformly on its whole support; for any proper selected set it
has the weaker upper bound $C_D\sqrt n$.
If $m\le cn$, it is bounded below by a positive constant on
$|\sum_I E_i-m\tau|\le R\sqrt n$ for
$n\ge n_0(c,D,R)\ge N_\Gamma$,
with constants depending on $c,D,R$ and the compact gap-mean range.
This threshold ensures that the displayed central set lies inside
the support of the remaining Gamma density.
In the exact-mean case $A=n\tau$, both upper bounds hold for every
$n\ge1$. The central lower bound still requires its stated threshold.
No uniform upper bound is asserted for a near-central bounded-$n$
window that permits $A\downarrow0$.
\end{lemma}
\begin{proof}
Independence and convolution of Gamma densities give \eqref{pr:RN}.
Stirling's formula, uniformly for $\tau$ in its compact interval, gives
\[
 \sup_x f_k(x)\le C k^{-1/2},\qquad
 f_k(k\tau+z\sqrt k)\asymp_R k^{-1/2}
 \quad(|z|\le R,\ k\ge N(R,\tau_-)).
\]
For $k=1$ the supremum is $1/\tau$ and the same upper bound holds.
For $n\ge N_\Gamma$, enlarge the threshold for the preceding central
estimate with $R=D$; the denominator in \eqref{pr:RN} is then comparable
to $n^{-1/2}$, and $A\ge n\tau/2$.
If $A=n\tau$, the explicit density and Stirling's bounds instead give
$f_n(n\tau)\asymp n^{-1/2}$ for every $n\ge1$, uniformly in $\tau$.
The asserted bounds follow; on the central selected set the remaining
total has displacement $O(\sqrt n)$ from a mean with shape comparable
to $n$. This is one joint density for all selected coordinates, not
independence of selected blocks after conditioning.
\end{proof}

\begin{proposition}[One-end exact-total bridge]\label{pr:gamma-one}
For fixed $D,D'$, the compact laws, and a fixed curve as below,
there are $N_{\rm one},H_0$
such that the following holds for $n\ge N_{\rm one}$.
Suppose $|\omega-\tau|\sqrt n\le D$, $|A-n\tau|\le D\sqrt n$, and
$H_0\le h\le D'\sqrt n$. The flat conditional survival probability has
$\lambda$ moment at most $C B_n(h)$, with the exponent of
\eqref{pr:ONE}. A stronger event has moment at least $c B_n(h)$ and can
retain a fixed curve $g\min(j,n-j)^\gamma$, $0<\gamma<1/6$, positive
unreserved terminal risk in $[a_0\sqrt n,a_1\sqrt n]$, and central
full-color windows. The constants and windows can be fixed uniformly
on the displayed parameter range.
One may enlarge $H_0$ after fixing $N_{\rm one}$ so that
$H_0>D'\sqrt{N_{\rm one}}$; the reserve range then itself enforces
the length threshold. The flat upper bound alone also holds for bounded
$n$, by increasing its constant.
\end{proposition}
\begin{proof}
All density comparisons below are made after taking
$N_{\rm one}\ge N_\Gamma$ and the required central-support thresholds.
Keeping just the first $\lfloor n/4\rfloor$ survival constraints and
using Lemma~\ref{pr:gamma} proves the upper bound. For the lower bound
first assume $h\le\varepsilon\sqrt n$. On $m=\lfloor n/4\rfloor$
slots retain the strong reference family in Theorem~\ref{pr:one}, with
\begin{equation}
 a_0\sqrt n\le S_m\le a_1\sqrt n,\quad
 |T_m-m\tau|\le R\sqrt n,\quad
 |N_c(m)-\pi_cm|\le R\sqrt n\quad\hbox{for every color }c.
 \label{pr:prefixboxes}
\end{equation}
Its conditional mark probability $f$ has $\mathbb E f^\lambda\ge cB_n(h)$.
The selected-gap density has a positive lower bound on these boxes.
Fix one retained prefix environment and one retained prefix label word.
For $k=n-m$, the remaining exact gap total $A-T_m$ is central. Represent
these gaps as
\[
 E'_i=(A-T_m)\frac{W_i}{W_1+\cdots+W_k},
 \qquad W_i\hbox{ iid exponential of mean one}.
\]
The centered polygonal partial gap sums, divided by $\sqrt k$, converge
to $\tau$ times a Brownian bridge plus their prescribed central linear
drift. Independently, centered color counts converge to multinomial
Brownian motion on the zero-sum hyperplane. Its weight projection has
positive variance. Both convergences are uniform over the central boxes
and compact laws: every parameter sequence has a convergent subsequence;
characteristic-function expansion identifies its Gaussian finite-dimensional
limits, and the uniform fourth-moment bound
$\mathbb E|\sum_{i=1}^k Z_i|^4\le C(k^2+k)$ gives tightness, including
polygonal intervals shorter than $1/k$. Normalizing by $k^{-1}\sum W_i$
preserves convergence by the law of large numbers.

The starting risk is in a fixed positive compact interval after scaling.
Choose a continuous color path in the zero-sum hyperplane whose weight
projection carries it to a fixed positive final window while remaining
strictly positive; choose the zero centered gap bridge. Each sufficiently
small tube about these paths has positive limiting probability. The
multinomial covariance is nondegenerate on its hyperplane, so all the
required color paths lie in its support. A finite set of tubes covers
the compact possible prefix endpoints and remaining central drifts.
The minimum probability is positive, and uniform convergence transfers
a fixed fraction to the finite bridges. Enlarging the final color windows
once absorbs the prefix deviations. On the tube the risk is at least
$c\sqrt n$ on the continuation, which dominates $gn^\gamma$ for large
$n$. On the first quarter the imposed one-sided curve is already the
required global curve. Thus every retained prefix word has joint
remaining-environment-and-mark continuation probability at least $c>0$.

For each prefix environment let $Z$ be the full successful mark probability
as a function of the remaining environment. Exactly
\begin{equation}
 0\le Z\le f,\qquad \mathbb E_{\rm remainder}Z\ge cf,
 \qquad \mathbb E_{\rm remainder}Z^\lambda\ge cf^\lambda.
 \label{pr:concavecontinue}
\end{equation}
The last inequality follows from $z^\lambda\ge f^{\lambda-1}z$ for
$0\le z\le f$. Integrate against the central lower density in
Lemma~\ref{pr:gamma}. This proves the lower bound. When $h$ is comparable
to $\sqrt n$, a fixed positive tube directly gives a constant lower
bound; alternatively use a smaller fixed multiple of $\sqrt n$ and
monotonicity. This is an annealed continuation for each prefix word,
not a quenched lower bound for every remaining environment. Central
full-type windows are retained; no individual exact type is asserted
to have constant probability.
The length threshold also includes the uniform tube-convergence and
endpoint-support thresholds after the windows are fixed. Indeed
$l-\omega=bp$ is bounded below by some $\beta>0$, whereas
$A\le n\omega+2D\sqrt n$. Thus
$nl-A\ge\beta n-2D\sqrt n$, which exceeds $a_1\sqrt n$
for sufficiently large $n$. For the opposite sign the corresponding
available endpoint is $A-n(l-b)\ge b(1-p)n-2D\sqrt n$, and $b(1-p)$
has the same kind of positive compact lower bound. Thus a positive
terminal window is not demanded beyond the possible mark total in
either orientation. The retained central color tubes above
give its actual support. No lower bound is obtained by compactness at
excluded small $n$. For the flat upper bound in the bounded range, use
$q\le1$ and $B_n(h)\ge n^{-1/2}$ for $h\ge1$.
\end{proof}

Fix a small $\epsilon>0$, $0<\gamma<1/6$, and put $d_n(i)=\min(i,n-i)$.
Let $\mathrm{REG}_K$ consist of the following inequalities for every
$0\le i<j\le n$ and every color $c$:
\begin{align}
 |T_j-T_i-\tau(j-i)|&\le\epsilon(j-i)
               +K[1+d_n(i)^\gamma+d_n(j)^\gamma],\notag\\
 |N_c(i,j)-\pi_c(j-i)|&\le\epsilon(j-i)
               +K[1+d_n(i)^\gamma+d_n(j)^\gamma].           \label{pr:REG}
\end{align}

\begin{lemma}[Relative regularity under the exact total]\label{pr:REGlemma}
Assume $n\ge1$, $A>0$, $|A-n\tau|\le D\sqrt n$,
$|\omega-\tau|\sqrt n\le D$, and $K\ge0$.
If $Q_{\rm bad}(E)$ is the conditional mark probability of flat survival
and $\mathrm{REG}_K^c$, then
\begin{equation}
 \mathbb E_A Q_{\rm bad}^\lambda\le C e^{-cK} B_n(h),
 \qquad h\ge1.                                             \label{pr:relativeREG}
\end{equation}
The constants are uniform in the stated near-centered regime.
\end{lemma}
\begin{proof}
Choose a fixed $N_{\rm REG}\ge\max\{32,N_\Gamma,(D/\epsilon)^2\}$,
enlarging it for the fixed-fraction comparisons below, and first take
$n\ge N_{\rm REG}$.
For a single bad interval $(i,j]$, exponential moments and Chernoff give
under the independent reference law
\begin{equation}
 \mathbb P(\mathcal D_{i,j})
 \le C\exp\{-c\epsilon^2(j-i)
          -c\epsilon K[1+d_n(i)^\gamma+d_n(j)^\gamma]\}.
 \label{pr:intervaltail}
\end{equation}
This follows by choosing a fixed tilt $z$ with
$\log\mathbb E e^{z(E-\tau)}\le Cz^2$ and $C|z|\le\epsilon/2$;
the finite color variables are even bounded.
We give a partition that preserves the survival scale under conditioning.

For $j\le n/4$, retain only survival to $n/2$. The event and the bad
interval use at most half the gaps. Lemma~\ref{pr:gamma} bounds their
conditional moment by the independent reference moment, to which
Lemma~\ref{pr:relativebad} applies. Its upper bound is
$C(1+j)^2\mathbb P(\mathcal D_{i,j})^{\lambda/4}B_n(h)$.
For $i\ge n/8$, retain necessary survival on the first $\lfloor n/16\rfloor$
slots. It is disjoint from the bad interval; their union omits a fixed
fraction of the gaps. After one joint density comparison they are
independent under the reference law, giving
$C\mathbb P(\mathcal D_{i,j})^\lambda B_n(h)$.
The remaining intervals have $i<n/8$, $j>n/4$ and length at least $n/8$.
For a proper gap interval use the $C\sqrt n$ density bound in
Lemma~\ref{pr:gamma}; for the entire gap set its total is fixed and its
linear tolerance is satisfied for large $n$. Mark distributions are
unaffected by the gap total. Thus all these remaining errors are
exponentially small, even after their $\lambda$ powers.

In the early region, the polynomial in $j$ is summable against
$e^{-cKj^\gamma}$ and the interval-length factor. In the late region
near the terminal end sum in $n-i$; away from both ends the additional
$e^{-cKn^\gamma}$ absorbs the number of intervals. There are at most
$n^2$ macroscopic intervals, and their exponential errors are smaller
than $e^{-cK}B_n(h)$ since $B_n(h)\ge c n^{-1/2}$.
Subadditivity of the $\lambda$ power proves \eqref{pr:relativeREG}.
For $1\le n<N_{\rm REG}$, every gap sub-sum is at most
$A\le N_{\rm REG}\tau_++D\sqrt{N_{\rm REG}}$, and every color count
is at most $N_{\rm REG}$. Consequently all of \eqref{pr:REG} hold
deterministically whenever
\[
 K\ge K_{\rm fin}:=1+2N_{\rm REG}\tau_+
                       +D\sqrt{N_{\rm REG}}+N_{\rm REG}.
\]
For $0\le K<K_{\rm fin}$ use $Q_{\rm bad}\le1$ and
$B_n(h)\ge N_{\rm REG}^{-1/2}$; increasing $C$ to at least
$e^{cK_{\rm fin}}\sqrt{N_{\rm REG}}$ proves the same estimate.
This handles the entire bounded-$n$ range even when $A\downarrow0$,
without a Gamma density comparison.
This proof does not subtract an unconditional error larger than the
survival moment.
\end{proof}

\subsection{Complete cells, every internal ordering, and Harris's inequality}

We specify the deterministic projection because its converse is needed
when a relative bad event is removed. Take $r$ independent forward
positions $U_i$ uniform on $[0,v]$, with $v$ an integer. Fix an integer
$N\ge1$, and set
\begin{equation}
 J_i=N+\left\lfloor\frac{v-N}{v}U_i\right\rfloor,
 \qquad J_i\in\{N,\ldots,v-1\}.                            \label{pr:cellmap}
\end{equation}
These are independent uniform retained cells and
$U_i-1<J_i\le U_i+N$. Let the slope per cell be $c>0$. With reverse
coordinates $\tau_i=v-U_i$, $\xi_i=v-J_i$, use strict before-jump deficits
\[
 R_{\rm full}(t)=ct-\sum_{\tau_i<t}w(C_i),\qquad
 R_{\rm cell}(t)=ct-\sum_{\xi_i<t}w(C_i).
\]
The displacement $\tau_i-N\le\xi_i<\tau_i+1$ and positivity of all
weights give exactly
\begin{equation}
 R_{\rm cell}(t)\le c+R_{\rm full}(t-1),\qquad
 R_{\rm full}(t)\le cN+R_{\rm cell}(t-N).                   \label{pr:twocell}
\end{equation}
The inequalities apply respectively for $t\ge1$ and $t\ge N$; earlier
times have bounds $c$ and $cN$. Integer before-jump tests control the
entire cell path since its deficit increases linearly between jumps.
In particular a cell flat cap implies a continuous flat cap with only
$cN$ extra reserve.

There are $n=r+1$ spacings between the $r$ genuine positions and both
endpoints. In physical units their law is that of $n$ exponential gaps
conditioned on total $cv$. Introduce one independent dummy mark at the
last spacing. Auxiliary after-jump and genuine before-jump risks differ
by at most $l$ at each tested time; removing the dummy changes each
color total by at most one. Fixed reserve and central-window adjustments
therefore transfer Proposition~\ref{pr:gamma-one} to these actual positions.
This transfer uses $n\ge N_{\rm one}$ after those adjustments;
it makes no endpoint-support assertion for excluded short clocks.
The last empty gap is retained throughout. This assertion concerns
central windows, not a prescribed exact type or a prescribed dummy color.

\begin{lemma}[Permutation-invariant regularity projection]\label{pr:cellREG}
There is a cell event depending only on the occupancy and the color
counts in each cell, implied by \eqref{pr:REG}, such that its occurrence
implies the required color interval bounds, with a larger fixed
coefficient, for \emph{every} ordering within cells. It also bounds each
cell occupancy by a fixed constant times one plus its distance to the
nearer physical endpoint raised to $\gamma$.
\end{lemma}
\begin{proof}
Here is a full deterministic formulation. Let ordered physical positions
$x_i=T_i$, $x_0=0,x_n=A$, be mapped to nondecreasing cells $J_i$ satisfying
$|aJ_i-x_i|\le\Delta_0$, with fixed $a,\Delta_0$; this includes the reverse
version of \eqref{pr:cellmap}. Put $D_A(t)=\min(t,A-t)$.
Applying the gap inequalities in \eqref{pr:REG} to $(0,i]$ and $(i,n]$
gives
\[
 (\tau-\epsilon)d_n(i)-K[1+d_n(i)^\gamma]
 \le D_A(x_i)\le
 (\tau+\epsilon)d_n(i)+K[1+d_n(i)^\gamma].
\]
The inequality $Kx^\gamma\le\eta x+C_{\eta,\gamma}K^{1/(1-\gamma)}$
therefore implies
\begin{equation}
 C^{-1}d_n(i)-B_K\le D_A(x_i)\le C d_n(i)+B_K.               \label{pr:rankdistance}
\end{equation}
The single-gap bound gives
$x_{i+1}-x_i\le\tau+\epsilon+K[1+d_n(i)^\gamma+d_n(i+1)^\gamma]$.
If a physical cut $t$ is within $\Delta_0$ of $[x_i,x_{i+1}]$, subtract
this bound from the lower bound in \eqref{pr:rankdistance} and absorb
its sublinear term. Thus
\begin{equation}
 d_n(i)^\gamma\le C D_A(t)^\gamma+B'_K.                     \label{pr:cutdistance}
\end{equation}
This includes $i=0,n-1$.

Let $I(s)=\#\{i:J_i\le s\}$, $g_s=I(s)-I(s-1)$,
$D(s)=D_A(as)$, and let $N_c(s,t)$ count color $c$ in cells $(s,t]$.
A cell cut $as$ is within $\Delta_0$ of $[x_{I(s)},x_{I(s)+1}]$.
Equations \eqref{pr:REG} and \eqref{pr:cutdistance} therefore imply
\begin{align}
 |N_c(s,t)-\pi_c[I(t)-I(s)]|
 &\le\epsilon[I(t)-I(s)]+K_1[1+D(s)^\gamma+D(t)^\gamma],\notag\\
 g_s&\le K_1[1+D(s)^\gamma].                               \label{pr:cell-envelopes}
\end{align}
For the second assertion, the first and last positions of an occupied
cell are separated by at most $a+2\Delta_0$; the gap lower bound then
bounds $g_s-1$, using \eqref{pr:cutdistance} at the adjacent cell cuts.
Furthermore, for each rank $u$ belonging to that cell,
\begin{equation}
 D(s)\le C d_r(u)+B''_K.                                   \label{pr:cell-rank}
\end{equation}
This follows from \eqref{pr:rankdistance} and the bounded displacement;
$|d_n(u)-d_r(u)|\le1$. It can be tested solely from occupancy counts,
since the smallest $d_r(u)$ in a cell's consecutive rank range occurs
at one of its rank endpoints.

Define the cell regularity event by \eqref{pr:cell-envelopes} and
\eqref{pr:cell-rank}. For an arbitrary rank interval in any internal
ordering, remove its two partial endpoint cells. Apply the first bound
of \eqref{pr:cell-envelopes} to the remaining complete-cell interval.
The two omitted parts have total error at most their occupancies, bounded
by the second assertion. Convert their physical endpoint distances back
to rank distances by \eqref{pr:cell-rank}. Adjacent cuts differ by a
fixed amount, and $|d_r(u+1)-d_r(u)|\le1$. Consequently, for every order,
\[
 |N_c(u,v)-\pi_c(v-u)|\le\epsilon(v-u)
          +K_2[1+d_r(u)^\gamma+d_r(v)^\gamma].
\]
If both endpoints are in one cell its entire contribution is controlled
by that cell's occupancy. The same argument covers the first and last
occupied cells and empty cuts. All constants are fixed before $n$ varies.
\end{proof}

\begin{theorem}[Critical complete-cell family]\label{pr:cell-critical}
Suppose $n=r+1\asymp T$, $|cv-n\tau|\le D\sqrt n$,
$|\omega-\tau|\sqrt n\le D$, and the retained-cell density is bounded
strictly below a fixed $\zeta<1$. Fix $0<\gamma<1/6$.
For $H_0\le h\le D'\sqrt T$ there is a family depending only on complete
colored cell occupancies whose conditional mark probability $f$ satisfies
\begin{equation}
 \mathbb E f^\lambda\ge c
       \min\{1,((1+h)/\sqrt T)^{2\chi(\lambda,\rho_\tau)}\}.
 \label{pr:cell-lower}
\end{equation}
The same family has the following properties: the ordinary forward
occupancy bounds $G_s\le s+1$ and $G_s\le\zeta s+C_0$; a reverse flat
cap with reserve $h+O(1)$; a prescribed fixed strong physical curve;
the complete-cell regularity of Lemma~\ref{pr:cellREG}, hence color
regularity for all internal orderings; a fixed positive central window
for the unreserved final risk; and central full-color windows.
The constants $C_0,N,H_0$ and all window widths are fixed. The assertion
uses $n\ge N_{\rm cell}$, with a fixed threshold chosen after these
data and containing the adjusted $N_{\rm one}$ and central-support
thresholds. If $n\ge c_*T$ is the fixed comparability bound, one may
finally enlarge $H_0>D'\sqrt{N_{\rm cell}/c_*}$ so that its reserve
range already enforces this length restriction.
More explicitly, the two reverse restrictions may be written
$R_{\rm cell}(k)\le h+C_1$ and
$R_{\rm cell}(k)\le h+C_g-g\min(k,v-k)^\gamma$ at every integer cut.
The flat extra reserve $C_1$ is independent of the selected amplitude $g$.

If a necessary competitor event skips at most $J_d\le C(1+d)^a$ slots
and retains a deterministic free block of a fixed positive fraction of
$n$, of length at most $c_1n$ with $c_1<1$, its cell conditional probability
$q_d$ obeys
\begin{equation}
 \mathbb E q_d^\lambda\le C(1+d)^C
       \min\{1,((1+h)/\sqrt T)^{2\chi(\lambda,\rho_\tau)}\}.
 \label{pr:cell-competitor}
\end{equation}
When a skip is too large to leave that block, the same bound holds after
increasing the fixed polynomial exponent.
\end{theorem}
\begin{proof}
First the ordinary occupancy conditions have fixed positive probability.
For independent retained cells with density at most $\zeta_0<\zeta$,
a binomial exponential bound gives
$\mathbb P(G_s>\zeta s+C_0)\le C e^{-cC_0-cs}$.
Choose $C_0$ large. For $G_s\le s+1$, the first $N$ cells are empty,
and the same negative-drift bound summed over $s\ge N$ has arbitrarily
small total when fixed $N$ is large. After deleting these fixed cells
the density still has a uniform margin for all sufficiently large $v$.
Thus their intersection has probability $c_{\rm occ}>0$. Both events
are increasing when any labeled forward position is moved rightward.

The order of transferring the strong event is essential. Apply
Proposition~\ref{pr:gamma-one} first to a stronger rank curve. Subtract
\eqref{pr:relativeREG}; on the remaining family the deterministic
comparisons \eqref{pr:rankdistance}--\eqref{pr:cutdistance} convert this
curve to the desired physical curve. The single-gap bound controls
before-jump interpolation as well. Now drop REG and all proof-only
prefix boxes, retaining the physical strong and flat conditions and the
position-independent total-type and final-risk windows. This larger
event still has the same lower fractional moment and is monotone in
each independent original forward position, for every fixed mark word.
If its conditional mark probability is $Q(U)$, Harris's product-measure
inequality gives
\begin{equation}
 \mathbb E[\boldsymbol1_{\rm occ}Q(U)^\lambda]
     \ge\mathbb P({\rm occ})\mathbb E Q(U)^\lambda.           \label{pr:Harris}
\end{equation}
It is applied to independent labeled positions, not to the dependent
Dirichlet spacings. For one coordinate the inequality is the identity
$2\operatorname{Cov}(f,g)=\mathbb E[(f(Z)-f(Z'))(g(Z)-g(Z'))]\ge0$;
induction proves the product version. Total-type restrictions do not
disturb monotonicity because they do not depend on the positions.

Project by \eqref{pr:cellmap}. Equation \eqref{pr:twocell} and the
H\"older modulus of $\min(t,v-t)^\gamma$ give a cell strong event and a
separate cell flat cap. The latter has only the fixed extra reserve $c$,
independent of the subsequent curve amplitude. By the converse in
\eqref{pr:twocell}, a bad cell family lies, for every compatible latent
position vector, in continuous flat survival with $cN+O(1)$ extra
reserve and failure of a stronger continuous regularity envelope.
Lemma~\ref{pr:REGlemma} bounds its fractional mass. Subtract this second
REG loss and use Lemma~\ref{pr:cellREG}. Pointwise comparisons are taken
before fractional powers: if $Q_{\rm full}(U)\le Q_{\rm cell}(J(U))$,
then integration gives the lower direction immediately; the converse
bad-event containment gives the upper direction for the error. No
incorrect reversal of conditional Jensen is used.

Here is a noncircular choice of constants. Fix the retained-cell deletion,
density margin, and $c_{\rm occ}$. Fix central endpoint boxes and a small
fraction of their flat comparison constant. Choose regularity tolerances
and $K$ so that \eqref{pr:relativeREG} pays for both REG losses and the
occupancy factor. Then choose the rank-curve amplitude needed in the
cell projection. Choose the relative KMT accuracy and finally enlarge
$H_0$ to dominate its fixed initial error and make the normalized curve
cost small. This is possible because $\gamma<1/2$; its normalized
amplitude decreases as $H_0$ grows. Relative error estimates are uniform
for reserves at least one, so the earlier choices do not change.
Finally fix $N_{\rm cell}$ above all length thresholds, and, if desired,
make the final reserve enlargement stated in the theorem. This proves
\eqref{pr:cell-lower}.

For a competitor, discarding $J_d$ positive bounded jumps increases its
necessary reserve by at most $lJ_d$. At the same rank index, replacing
latent positions by cell centers changes only the two endpoint times
by a fixed amount, not by an amount proportional to the rank. Its free
block is therefore majorized by reference survival with reserve
$h+C+lJ_d$. Since $n\ge N_{\rm cell}\ge N_\Gamma$,
Lemma~\ref{pr:gamma} and Theorem~\ref{pr:one} give
\eqref{pr:cell-competitor}. If the skip exceeds a fixed fraction of $T$,
use $q_d\le1$; a larger fixed power of $1+d$ absorbs one since
$B_T(h)\ge cT^{-1/2}$. This completes the same-family assertion.
\end{proof}

In the critical application $h=H_0+M/4$, $0\le M\le\sqrt T$,
$\lambda=\theta+O_D(T^{-1/2})$, and
$\rho_\tau=\rho_c+O_D(T^{-1/2})$. Proposition~\ref{pr:lipschitz}
shows that the factors in \eqref{pr:cell-lower} and
\eqref{pr:cell-competitor} are uniformly comparable to
\[
 \min\{1,((1+M)/\sqrt T)^{2\chi(\theta,\rho_c)}\}.
\]
Indeed the error in its logarithm is $O_D(\log T/\sqrt T)$.
The assertions hold for every fixed central first count; mixing a fixed
central Poisson slab of positive probability preserves the lower bound.
Here the original length and each retained central count are taken above
the fixed $N_{\rm cell}$ threshold, as required by the lower construction.

\subsection{Positive drift: an effective prefix and uniform continuation}

Continue with \eqref{pr:walk}, now $0<\delta=\omega-\tau\le\delta_0$,
and condition on $T_n=n\tau$. Put
\begin{equation}
 B_{n,\delta}(h)=
 \min\{1,[(1+h)(\delta+n^{-1/2})]^{2\chi(\lambda,\rho_\tau)}\}.
 \label{pr:driftscale}
\end{equation}

\begin{theorem}[Positive-drift complete-cell family]\label{pr:cell-drift}
For fixed $\gamma\in(0,1/6)$ and fixed curve and occupancy requirements,
there are $\delta_0,H_0,D_0>0$ such that, for
$0<\delta\le\delta_0$, $\delta\sqrt n\ge D_0$, $h\ge H_0$,
the flat exact-total survival moment is at most $C B_{n,\delta}(h)$,
and a complete-cell family with the properties of
Theorem~\ref{pr:cell-critical} has moment at least $c B_{n,\delta}(h)$.
Its final unreserved risk can be restricted to
$[\delta n-C\sqrt n,\delta n+C\sqrt n]$, and its full color counts are
central. Deterministic competitor skips satisfy the corresponding upper
bound $C(1+J)^C B_{n,\delta}(h)$.
\end{theorem}

\begin{proof}
Fix constants $A_0,K_0$, and set $m=\lfloor A_0/\delta^2\rfloor$.
Choose $D_0$ later so $m\le n/8$, and decrease $\delta_0$ to absorb
rounding. On this effective prefix $\delta\sqrt m$ is bounded. Apply
Theorem~\ref{pr:one}, retaining
\begin{equation}
 S_m\in[K_0/\delta,C_{A_0,K_0}/\delta],\quad
 |T_m-m\tau|\le C_{A_0,K_0}/\delta,\quad
 |N_c(m)-\pi_cm|\le C_{A_0,K_0}/\delta.                     \label{pr:driftprefix}
\end{equation}
Its conditional probability $f$ has moment at least
$c\min(1,[(1+h)\delta]^{2\chi})$. For $h\delta$ bounded below, a fixed
tube at the effective scale gives this constant lower bound directly;
large reserves help. Lemma~\ref{pr:gamma} bounds the prefix density
above and below on \eqref{pr:driftprefix}. The upper survival bound
already follows by retaining only this prefix.
For this application the complete total is exactly $n\tau$, so the
upper comparison uses the all-$n$ exact-mean clause. The central lower
comparison uses its length threshold, enforced by the eventual choices
of large $D_0$ and small $\delta_0$.

Fix its environment and a retained label word. Let $k=n-m$ and
$R=n\tau-T_m$. Retain remaining full-color totals within fixed
$C_1\sqrt k$ windows around $k\pi_c$. Their joint multinomial probability
is bounded below. Conditional on any such type, use independent reference
parameters equal to its proportions $\pi'_c$ and to $\tau'=R/k$.
The binary sequence is the uniform fixed-count permutation and the gaps
are a Dirichlet bridge. The new parameters satisfy
\[
 \pi'_c=\pi_c+O(C_1/\sqrt n),\qquad
 \tau'=\tau+O(C_{A_0,K_0}/(\delta n)),\qquad
 \delta':=\sum_c\pi'_cw(c)-\tau'\ge\delta/2,
\]
after $D_0$ is large in terms of the constants already fixed. Also
$\delta'\le C\delta$. The complete endpoint then belongs to the stated
$\delta n+O(\sqrt n)$ window and is at least $c\delta n$.

We verify the required continuation estimates under these two exact
conditions. A selected initial or terminal block of at most $k/2$ gaps
has density bounded by a constant relative to the new reference law.
Indeed its full total is exactly $R=k\tau'$ with $\tau'$ in a positive
compact interval, so this is again the all-$k$ exact-mean upper bound.
The corresponding binary density ratio is
\[
 \frac{\mathbb P\{\operatorname{Bin}(k-j,p')=u-U_j\}}
      {\mathbb P\{\operatorname{Bin}(k,p')=u\}},\qquad u=kp'.
\]
Its denominator is comparable to $k^{-1/2}$ and numerator at most
$C(k-j)^{-1/2}$, so it too is bounded. The two ratios multiply by
independence. Specifying all colors leaves this same binary bridge law.

For independent reference increments $X=w(C)-E$, uniform exponential
moments and Taylor expansion give, for a fixed small $a>0$,
\[
 \mathbb E\exp[-a\delta(X-\delta'/2)]\le1.
\]
The exponential supermartingale therefore gives, from $x\ge K_0/\delta$,
\begin{equation}
 \mathbb P\{\exists j\ge0:x+\textstyle\sum_{i\le j}X_i
                      <x/2+\delta'j/2\}\le e^{-c\delta x}.
 \label{pr:ruin}
\end{equation}
The bounded initial-half density transfers this estimate to the first
half of the conditioned continuation. For its late half, the maximal
Bernstein estimate for each independent centered channel is
\[
 \mathbb P\{\max_{j\le k/2}|\textstyle\sum_{i\le j}Z_i|>z\}
          \le C\exp[-c\min(z^2/k,z)].
\]
Apply the same bounded density to the initial half and to the reversed
terminal half. Both centered totals are fixed to zero; these halves
cover the complete bridge, giving
\begin{equation}
 \mathbb P\{\max_{j\le k}|S'_j-\delta'j|>c_2\delta k\}
                    \le C e^{-c\delta^2 k}.                \label{pr:bridge-max}
\end{equation}
There is no factor $k$ from a union over individual times.

The global strong curve on the continuation is at most $g(m+j)^\gamma$.
At its beginning $gm^\gamma=o(\delta^{-1})$, and its derivative is
$O(\delta^{2-2\gamma})=o(\delta)$ as $\delta_0\to0$.
Thus the positive linear margin in \eqref{pr:ruin} dominates it on the
first half. On the late half $gn^\gamma=o(\delta n)$ uniformly when
$\delta\sqrt n\ge D_0$, so \eqref{pr:bridge-max} suffices.
Choose $K_0$ so $Ce^{-cK_0}<1/4$, then $D_0$ so the late failure is
less than $1/4$ and all parameter comparisons hold. Every retained
prefix word and remaining type now has conditional continuation
probability at least $1/2$, averaged over its remaining exact gaps.
The central type window has fixed positive mass. Thus the inequalities
\eqref{pr:concavecontinue} hold again, proving the strong exact-total
lower bound after integrating the central prefix density.

For completeness, regularity is not inferred from endpoint centrality.
Before exact conditioning, the same effective-prefix argument and
\eqref{pr:ruin} prove flat reference survival comparable to
\eqref{pr:driftscale}: use the near-centered theorem when
$\delta\sqrt n$ is bounded, and the effective prefix otherwise.
Hence time doubling and polynomial reserve comparison hold uniformly
throughout this positive-drift layer. Repeat the three-region proof of
Lemma~\ref{pr:REGlemma} with the unchanged exponential interval tails.
It gives a relative bad moment at most
$Ce^{-cK}B_{n,\delta}(h)$ under the exact total. The ordinary occupancy,
physical projection, and two-stage REG/Harris argument are now precisely
those in Theorem~\ref{pr:cell-critical}. The proof-only effective-prefix
boxes are dropped before Harris. The final endpoint and total-type
windows are position independent and remain.

Constants can be fixed in the following order: occupancy/deletion
requirements; effective-prefix constants $A_0,K_0$ and their endpoint
fraction; central type widths; REG tolerance and $K$; required curve
amplitude and relative coupling accuracy; minimum reserve $H_0$;
then large $D_0$, small $\delta_0$, and the length threshold. The
reference relative errors are uniform in $h\ge1$. Enlarging $H_0$ and
shrinking $\delta_0$ make normalized curve costs small without changing
the earlier fixed success fraction. The two displayed parameter
inequalities enforce $n\ge(D_0/\delta_0)^2$; take these final choices
so that this exceeds every required length threshold. Competitor blocks
have enlarged
reserve $h+C+lJ$ and bounded selected-gap density. Use the reference
upper bound; if $J$ is a fixed fraction of the effective time
$\min(n,\delta^{-2})$, the trivial bound one is absorbed by a larger
polynomial. This proves all assertions.
\end{proof}

In the ordinary boundary application, $n\asymp T$ and the count-specific
drift equals $d/\kappa+O(T^{-1/2})$, with $\kappa$ bounded above and below.
The reference parameters differ from their critical values by
$O(d+T^{-1/2})$. Proposition~\ref{pr:lipschitz} permits the critical
exponent in the scale
\[
 \min\{1,[(1+h)(d+T^{-1/2})]^{2\chi_c}\}.
\]
Indeed $d\log(1/d)$ and $T^{-1/2}\log T$ are bounded. The region of
bounded $d\sqrt T$ is already covered by
Theorem~\ref{pr:cell-critical}; no intermediate drift range is omitted.

\subsection{Physical Poisson time and weighted insertion}

Let $\Pi$ be a homogeneous Poisson process of rate $\kappa$ on $[0,s]$,
where $\kappa$ ranges over a fixed positive compact interval. Its points
carry independent marks of the preceding law. For fixed $c>0$ write
\[
 S(t)=\sum_{u\in\Pi,\ u\le t}w(C_u)-ct,
 \qquad \tau=c/\kappa,\qquad
 \rho_\tau=\frac{\tau^2}{\tau^2+\operatorname{Var}(w)}.
\]
Before-jump and after-jump conventions differ by at most $l$ and will
always be allowed that fixed reserve adjustment.

\begin{theorem}[Physical one-end upper bound]\label{pr:physical-one}
Fix $0<\gamma<1/2$, $g\ge0$, $H\ge1$. Given $\Pi$, let $q$ be the
conditional mark probability that either $S(t)$ or $-S(t)$ stays above
$-H-gt^\gamma$ for every $0\le t\le s$. If
$|\kappa\omega-c|\sqrt s\le D$, then
\begin{equation}
 \mathbb E q^\lambda\le C_{D,g,\gamma}
       \min\{1,[(1+H)/\sqrt s]^{2\chi(\lambda,\rho_\tau)}\}.
 \label{pr:physical-upper}
\end{equation}
The correlation is the arrival-gap correlation $\rho_\tau$ displayed
above. The bound also permits a fixed logarithmic wall in place of
the power wall.
\end{theorem}
\begin{proof}
We first note why the soft-wall form of Theorem~\ref{pr:one} has its
stated uniformity. After a bad prefix of length $b$,
$(b+j)^\gamma\le b^\gamma+j^\gamma$ changes reserve by only a fixed
polynomial in $b$. Thus the relative bad-prefix proof is unchanged up
to its fixed polynomial. Under Lamperti transformation from time $H^2$,
the wall becomes $gH^{2\gamma-1}e^{-(1/2-\gamma)u}$, whose action and
coefficient variation are bounded for $H\ge1$. The relative KMT proof
and its absorption therefore hold for both signs; no fixed unconditional
REG error is subtracted to prove this upper bound.

Take $m=\lfloor\kappa s/2\rfloor$; bounded $s$ is absorbed in constants.
The Poisson lower-tail probability $\mathbb P(N_s<m)$ is at most
$e^{-c_0s}$. Extend the process beyond $s$ if necessary. Its first $m$
scaled arrival gaps $E_i=c(t_i-t_{i-1})$ are independent exponentials
of mean $\tau$. The extension does not alter the conditional mark
probability on $[0,s]$.

On $N_s\ge m$, physical survival imposes necessary inequalities at these
first $m$ arrivals. For the sign $S$, at a relevant test time the
accumulated weight is at most $lj$; hence survival implies
$ct\le lj+H+gt^\gamma$. Since $\gamma<1$, it follows that
$t\le C(j+H+1)$ and consequently
$gt^\gamma\le C(j^\gamma+H+1)$. Thus the necessary rank event has reserve
$C(H+1)$ and a fixed soft-power coefficient. For the sign $-S$, no
arrival-time upper bound is needed when $-S(t)\ge0$: its rank lower
wall is automatic. When $-S(t)<0$, one has $ct<\sum w\le lj$, which
again bounds the allowed physical soft wall by a fixed rank power wall.
This separate argument avoids reversing an invalid time inequality.

At a fixed full environment the mark probability of this necessary rank
event depends only on the first $m$ gaps. Dropping the restriction
$N_s\ge m$ in its expectation gives
\[
 \mathbb E q^\lambda\le\mathbb E q_{\rm rank}^\lambda+e^{-c_0s}.
\]
Here $|\omega-\tau|\sqrt m\le C_D$. Apply Theorem~\ref{pr:one} with
its soft wall to the iid arrival gaps. Since $m\asymp s$ this gives
\eqref{pr:physical-upper}. Finally the displayed scale is at least
$cs^{-1/2}$, because $\chi\le1/2$ and $H\ge1$, so the exponential tail
is absorbed. A logarithm is bounded by a fixed constant plus any fixed
positive power $t^\gamma$, proving the last assertion.
\end{proof}

\begin{corollary}[Effective physical time]\label{pr:physical-drift}
If the physical mean $d=\kappa\omega-c\ge0$ is sufficiently small, the
same full-horizon necessary event has fractional moment at most
\[
 C\min\{1,[(1+H)(d+s^{-1/2})]^{2\chi(\lambda,\rho_\tau)}\}.
\]
The bound is valid for either risk sign and the preceding soft walls.
\end{corollary}
\begin{proof}
Retain only the deterministic interval of length
$s_0=c_0\min(s,d^{-2})$, with $d^{-2}=\infty$ when $d=0$.
Then $d\sqrt{s_0}$ is bounded. Apply Theorem~\ref{pr:physical-one};
$s_0^{-1/2}\asymp d+s^{-1/2}$. When $s_0$ is bounded the conclusion
follows from the trivial upper bound one. Restricting to this interval
is a necessary condition for full survival, so no continuation lower
bound is asserted or required.
\end{proof}

\begin{proposition}[Palm bounds with the survival weight retained]
\label{pr:palm}
Let $T\ge1$ and let $q_H(\Pi)$ be a physical conditional survival
probability of the preceding kind on $[0,T]$. Suppose its proved upper
bound is $C B_T(H)$, where either
\[
 B_T(H)=\min\{1,[(1+H)/\sqrt T]^{2\chi}\}
 \quad\hbox{or}\quad
 B_T(H)=\min\{1,[(1+H)(d+T^{-1/2})]^{2\chi}\}.
\]
For $t\in[0,T]$ set
\[
 Z(t)=1+\sum_{u\in\Pi,\ u\ge t}e^{-(u-t)},\qquad
 \mathcal A=1+C_0\sup_{0\le t\le T}\frac{Z(t)^5}{(2+t)^4},
 \qquad N_T=|\Pi|.
\]
For each integer $m\ge0$ and every fixed $a>1/4$,
\begin{align}
 \mathbb E[Z(t)^m q_H^\lambda]&\le C_m B_T(H),\notag\\
 \mathbb E[N_T Z(t)^m q_H^\lambda]&\le C_m T B_T(H),\notag\\
 \mathbb E[(T+N_T)\mathcal A^a q_H^\lambda]
                                      &\le C_a T B_T(H).    \label{pr:palmweighted}
\end{align}
Moreover for $k\ge0$,
\begin{equation}
 \mathbb E[(T+N_T)
       q_{H+k+\lceil\log\mathcal A\rceil}^{\lambda}]
                          \le C(1+k)^C T B_T(H).             \label{pr:randomreserve}
\end{equation}
\end{proposition}
\begin{proof}
Adding $j$ deterministic points changes every risk by at most $jl$,
regardless of their marks. Integrating the new marks therefore gives
the pointwise inequality
\begin{equation}
 q_H(\Pi+\{u_1,\ldots,u_j\})\le q_{H+jl}(\Pi).             \label{pr:insert}
\end{equation}
It holds for either sign. The proved survival upper bounds imply
$\mathbb E q_{H+jl}^\lambda\le C(1+j)^C B_T(H)$, since $H\ge1$ and
$2\chi\le1$.

Expand $Z(t)^m$ and group monomials by their distinct Poisson points.
The multivariate Mecke identity for $j$ distinct points says that
\[
 \mathbb E\sum_{u_1,\ldots,u_j\in\Pi}^{\ne}
 G(u_1,\ldots,u_j,\Pi)
 =\kappa^j\int_{[0,T]^j}
   \mathbb E G(u_1,\ldots,u_j,\Pi+\{u_1,\ldots,u_j\}),d\boldsymbol u.
\]
For a Poisson process this follows directly by expanding its factorial
probabilities $e^{-\kappa T}(\kappa T)^n/n!$, choosing the $j$ ordered
distinct points, and reindexing $n-j$; thus no independence of the weight
and the environment is being assumed. In each monomial the insertion
integrals have bounded mass, since every point carries at least one
factor $e^{-(u-t)}\boldsymbol1_{\{u\ge t\}}$. Use
\eqref{pr:insert} after insertion to obtain the first bound in
\eqref{pr:palmweighted}. One additional Mecke point without a decaying
factor gives a factor $\kappa T$ and proves the second; coincidences
with existing monomial points only reduce the number of integrations.

For $j\le t<j+1$ one has $Z(t)\le eZ(j)$. Hence
$\mathcal A^a\le C_a[1+\sum_{j\ge0}(2+j)^{-4a}Z(j)^{5a}]$.
Bound the noninteger moment by a larger integer moment and sum the
convergent series, since $4a>1$. This proves the last assertion of
\eqref{pr:palmweighted}.
Finally, on $\lceil\log\mathcal A\rceil=j$ one has
$\mathcal A^a\ge e^{a(j-1)}$. Apply the weighted estimate with reserve
$H+k+j$, separately in each bin, to get
\[
 \mathbb E[(T+N_T)q_{H+k+\lceil\log\mathcal A\rceil}^{\lambda}]
 \le CT\sum_{j\ge0}e^{-a(j-1)} B_T(H+k+j)
 \le C(1+k)^C T B_T(H).
\]
The reserve in the weighted estimate depends on the bin $j$; a single
unweighted crowding tail would not justify this step. This proof also
shows why arbitrary stopped depths may be controlled by the global
crowding majorant without claiming that a stop remains fixed when
unseen Poisson points are deleted.
\end{proof}

The probability results proved here are inputs about conditional mark
families. Their use for divisor-box volumes requires the separate positive
complete-cell kernel, geometric cover, and terminal-count arguments.
In particular no pointwise identification of an actual divisor union
with a fractional survival probability is used in any proof in this section.
